This chapter details how the proposed testbed and capabilities are implemented at \acs{ITAv}. The architecture of the implemented testbed is presented in Section~\ref{chap:impl:architecture}, identifying the concrete components responsible for each of the proposed capabilities and how they interact. The subsequent sections detail the implementation of each capability, following the order adopted in the proposal:~\acs{5G}~resource lifecycle management in Section~\ref{chap:impl:resource_lifecycle}, network exposure in Section~\ref{chap:impl:exposure}, network observability in Section~\ref{chap:impl:observability}, and service delivery in Section~\ref{chap:impl:delivery}. Section~\ref{chap:impl:capabilities} then summarizes the full set of capabilities offered by the testbed.

Finally, the multi-vertical use case defined in the proposal is implemented in Section~\ref{chap:impl:use-case}, demonstrating how verticals and the network application can interact with the \acs{5G} system through the exposed interfaces. This includes the provisioning of the required \acs{5G} resources, the enforcement of \acs{QoS} requirements, and the monitoring of network conditions to support application and network \acs{QoS} adaptation.
